\section{Exact Non-Forgetting, Operational Memory}
\label{sec:nonforget}

The term \emph{without forgetting} has a precise scope in this paper:
the predictor channel is non-dissipative on the joint density matrix,
even though individual downstream probes may not retain their accuracy
at every horizon. For $P_k(\rho)=U^k\rho(U^\dagger)^k$, unitary
conjugation preserves the joint spectrum and therefore every spectral
functional of the joint state~\cite{nielsen2010quantum}:
\begin{equation}
  \begin{aligned}
  \mathrm{spec}(P_k(\rho))&=\mathrm{spec}(\rho),\\
  \Tr(P_k(\rho)^2)&=\Tr(\rho^2),\\
  S(P_k(\rho))&=S(\rho).
  \end{aligned}
  \label{eq:joint-invariants}
\end{equation}
These identities encode the structural reason the predictor cannot
dissipate information during blind
rollout~\cite{nielsen2010quantum,baker2022quantum}.
The implementation preserves these invariants to numerical precision
over $30$ rollout steps; the largest observed drift is below
$2.4\times 10^{-7}$ (Supplementary \cref{fig:supp-invariants}).

The qualification matters as much as the guarantee. The JEPA loss and
the probes observe $\Tr_{\env}\rho$ through a projection rather than
the full joint state, and partial trace is not unitary. Information can
move into system--environment correlations and become inaccessible to a
reduced-state linear probe while the joint state remains exactly
non-dissipative. The paper therefore reports two levels of evidence:
\cref{eq:joint-invariants} fixes the structural property of the
predictor, and the blind-rollout experiments in \cref{sec:imagine}
measure how much task-relevant information survives reduction,
projection, and repeated prediction without new observations.

The same geometry yields a useful diagnostic for what a unitary
predictor can match on the full joint space. The spectrum-mismatch
theorem in \cref{sec:theorem} states that
\begin{equation}
  \min_{U\in \mathrm{U}(\dtot)}
  \|U\rho U^\dagger-\sigma\|_F^2
  =
  \|\lambda^\downarrow(\rho)-\lambda^\downarrow(\sigma)\|_2^2.
  \label{eq:spectrum-mismatch-short}
\end{equation}
A unitary predictor can therefore match a full-state target exactly
only when context and target density matrices are isospectral. The
statement constrains the predictor's orbit geometry on the joint state
and does not transfer as a lower bound to the reduced and projected
JEPA loss in \cref{eq:jepa-loss}. In training, the online and EMA
target encoders co-adapt, so the objective can become small by learning
spectrum-compatible regions of latent space before the
reduced/projected loss is evaluated. The full theorem, proof, and a
numerical orbit visualisation appear in the supplement.
